\subsection{Hermitian symmetric spaces}
We begin by recalling useful facts and definitions concerning Hermitian symmetric spaces and setting up our notation. Good references for what follows are \cite{jep93Wol72,jep93AMRT10,jep93Hel01,jep93Kna02}.

Let $G_{\mathbb R}$ be a simple real algebraic Hermitian Lie group and $K_{\mathbb R}$ a maximal compact subgroup of $G_{\mathbb R}$. By our convention, this means that $G_{\mathbb R}$ is a simple noncompact Lie group whose associated symmetric space $M=G_{\mathbb R}/K_{\mathbb R}$ is a Hermitian symmetric space of the noncompact type ($M$ is irreducible since $G_{\mathbb R}$ is simple), and that moreover $G_{\mathbb R}$ is the connected component of the group of real points $\jepPubBmi{G}(\mathbb R)$ of an algebraic group $\jepPubBmi{G}$ defined over $\mathbb R$. Observe that for example the connected component $\operatorname{Isom}^{\circ}(M)$ of the group of isometries of a Hermitian symmetric space of the noncompact type $M$ is a real algebraic Hermitian Lie group, see e.g. \cite[p.~106]{jep93AMRT10}. The group of complex points $\jepPubBmi{G}(\mathbb C)$ of $\jepPubBmi{G}$ will be denoted by $G$.

We shall first assume that the complex Lie group $G$ is simply connected. This assumption simplifies the exposition of the paper and the arguments of our proofs. We shall see in Remarks~5.14 and~5.15 how to deal with the general case.

The rank of the symmetric space $M$, or equivalently the real rank of $G_{\mathbb R}$, will be denoted by $r_M$. Recall that a $k$-dimensional \emph{flat} in $M$ is the image of a totally geodesic isometric embedding of $\mathbb R^k$ in $M$. By definition the rank $r_M$ of $M$ is the maximal dimension of a flat in $M$. Maximal flats are conjugated under $G_{\mathbb R}$ and maximal flats through the origin $eK_{\mathbb R}$ are conjugated under $K_{\mathbb R}$. Equivalently, in the Hermitian setting, one can consider \emph{polydiscs} in $M$, that is, images of totally geodesic homomorphic embeddings of a polydisc $\Delta^k$ in $M$. Then $r_M$ is also the maximal (complex) dimension of a polydisc in $M$ (maximal polydiscs are complexifications of maximal flats).

The fact that $M$ is an irreducible Hermitian symmetric space is equivalent to the fact that the Lie algebra $\mathfrak k_{\mathbb R}$ of $K_{\mathbb R}$ has a 1-dimensional center $\mathfrak z_{\mathbb R}$ and that the centralizer of $\mathfrak z_{\mathbb R}$ in the Lie algebra $\mathfrak g_{\mathbb R}$ of $G_{\mathbb R}$ is $\mathfrak k_{\mathbb R}$, see e.g. \cite{jep93Kna02}. If $\mathfrak g_{\mathbb R}=\mathfrak k_{\mathbb R}\oplus\mathfrak m_{\mathbb R}$ is a Cartan decomposition of $\mathfrak g_{\mathbb R}$, the complex structure at the origin is given by the adjoint action on $\mathfrak m_{\mathbb R}$ of a suitably normalized element of $\mathfrak z_{\mathbb R}$. It follows that any maximal Abelian subspace $\mathfrak t_{\mathbb R}$ of $\mathfrak k_{\mathbb R}$ contains $\mathfrak z_{\mathbb R}$ and is a Cartan subalgebra of $\mathfrak g_{\mathbb R}$. The corresponding subgroup $T_{\mathbb R}\subset K_{\mathbb R}\subset G_{\mathbb R}$ is a torus.

The Lie algebra $\mathfrak g$ of $G$ is the complexification of the Lie algebra $\mathfrak g_{\mathbb R}$ of $G_{\mathbb R}$. We denote by $\mathfrak z\subset\mathfrak t\subset\mathfrak k\subset\mathfrak g$ the complexifications of the subalgebras $\mathfrak z_{\mathbb R}\subset\mathfrak t_{\mathbb R}\subset\mathfrak k_{\mathbb R}$ and by $Z\subset T\subset K$ the corresponding complex algebraic subgroups of $G$. We also let $\mathfrak m$ be the complexification of $\mathfrak m_{\mathbb R}$, so that $\mathfrak g=\mathfrak k\oplus\mathfrak m$ is a Cartan decomposition of $\mathfrak g$. In particular, $[\mathfrak k,\mathfrak k]\subset\mathfrak k$, $[\mathfrak k,\mathfrak m]\subset\mathfrak m$ and $[\mathfrak m,\mathfrak m]\subset\mathfrak k$.

Let $z\neq0$ be an element in the center $\mathfrak z$ of $\mathfrak k$. Since the adjoint action of $\mathfrak z_{\mathbb R}$ gives the complex structure of $\mathfrak m_{\mathbb R}$, $\operatorname{ad}(z)\jepPubRestrict_{\mathfrak m}$ has exactly two opposite eigenvalues. We choose $z$ so that these eigenvalues are $2$ and $-2$ and we let
\[
\mathfrak m^+=\{x\in\mathfrak m\mid\operatorname{ad}(z)x=2x\},\qquad
\mathfrak m^-=\{x\in\mathfrak m\mid\operatorname{ad}(z)x=-2x\}
\]
be the corresponding eigenspaces. These are abelian subspaces of $\mathfrak m$.

This in turn implies that $\mathfrak p=\mathfrak k\oplus\mathfrak m^+$ is a maximal parabolic subalgebra of $\mathfrak g$. Let $P$ be the corresponding parabolic subgroup of $G$ and $\check M$ the projective variety $G/P$. Then $\check M$ is a symmetric space of the compact type called the compact dual of $M$. Indeed if we let $\check{\mathfrak g}=\mathfrak k_{\mathbb R}\oplus i\mathfrak m_{\mathbb R}$, $\check{\mathfrak g}$ is the Lie algebra of a real compact form $\check G_{\mathbb R}$ of $G$, and $\check M=\check G_{\mathbb R}/K_{\mathbb R}$. The symmetric space $M$ is then an open $G_{\mathbb R}$-orbit in $\check M=G/P$.

\begin{example}
For $p\leqslant q$ let $G_{\mathbb R}=\mathrm{SU}(p,q)$ be the special unitary group of the Hermitian form of signature $(p,q)$ on $\mathbb C^{p+q}$ whose matrix in the canonical basis is $I_{p,q}=\operatorname{diag}(-1,\ldots,-1,1,\ldots,1)$. Then one can choose $K_{\mathbb R}=\mathrm S(\mathrm U(p)\times\mathrm U(q))$ and the symmetric space $M$ identifies with the set of $p$-planes in $\mathbb C^{p+q}$ on which the Hermitian form is negative definite. It has rank $p$. As a bounded symmetric domain, it identifies with $\{Y\in \jepPubScript{M}_{p,q}(\mathbb C)\mid I_p-Y^{\star}Y\text{ is positive definite}\}$. We have $G=\mathrm{SL}(p+q,\mathbb C)$, $K=\mathrm S(\mathrm{GL}(p)\times\mathrm{GL}(q))$, $T$ is the torus of diagonal matrices in $G$, $Z$ is the 1-dimensional subgroup of diagonal matrices of the form $\operatorname{diag}(\lambda^q,\ldots,\lambda^q,\lambda^{-p},\ldots,\lambda^{-p})$ for $\lambda\in\mathbb C^{\star}$, $\mathfrak m\simeq\mathbb C^{pq}\times\mathbb C^{pq}$ is the subspace of block anti-diagonal matrices in $\mathfrak{sl}(p+q,\mathbb C)$. The compact real form of $G$ is $\mathrm{SU}(p+q)$ and the compact dual $\check M$ identifies with the Grassmannian of $p$-planes in $\mathbb C^{p+q}$, of which $M$ is an open subset.
\end{example}

Let $R$ be the set of roots of $G$. For $\alpha\in R$, $\mathfrak g_\alpha\subset\mathfrak g$ is the root space of $\alpha$, $\mathfrak{sl}(\alpha)$ is the Lie subalgebra of $\mathfrak g$ generated by $\mathfrak g_\alpha$ and $\mathfrak g_{-\alpha}$ and $\mathrm{SL}(\alpha)$ the corresponding subgroup of $G$.

A root space $\mathfrak g_\alpha$, $\alpha\in R$, is either a subspace of $\mathfrak k$ or a subspace of $\mathfrak m$, and the root $\alpha$ is said to be \emph{compact} or \emph{noncompact} accordingly. Equivalently, a root $\alpha$ is compact if and only if $\langle\alpha,z\rangle=0$. We have:
\[
\mathfrak k=\mathfrak t\oplus\jepPubOplus_{\alpha:\langle\alpha,z\rangle=0}\mathfrak g_\alpha,\quad
\mathfrak m=\jepPubOplus_{\alpha:\langle\alpha,z\rangle\neq0}\mathfrak g_\alpha,\quad
\mathfrak m^+=\jepPubOplus_{\alpha:\langle\alpha,z\rangle=2}\mathfrak g_\alpha,\quad
\mathfrak m^-=\jepPubOplus_{\alpha:\langle\alpha,z\rangle=-2}\mathfrak g_\alpha.
\]
We choose a basis $\Pi$ of $R$ so that the roots in $R(\mathfrak m^+)$ are positive roots. If $\alpha\in R$, we write $\alpha=\sum_{\beta\in\Pi}n_\beta(\alpha)\beta$ the expression of $\alpha$ in terms of the simple roots. The support $\operatorname{supp}(\alpha)$ of $\alpha\in R$ is the set $\{\beta\in\Pi\mid n_\beta(\alpha)\neq0\}$.

The Weyl group of $R$ is denoted by $W$. If $\alpha\in R$, $\alpha^\jepPubCoroot\in\mathfrak t$ denotes the corresponding coroot: it is defined by the relation $s_\alpha(\beta)=\beta-\langle\beta,\alpha^\jepPubCoroot\rangle\alpha$ for all $\beta\in\mathfrak t^*$, where $s_\alpha\in W$ is the reflexion corresponding to $\alpha$ (in particular, $s_\alpha(\alpha)=-\alpha$ and hence $\langle\alpha,\alpha^\jepPubCoroot\rangle=2$). If $\alpha=\sum_{\beta\in\Pi}n_\beta(\alpha)\beta$ is the expression of a root $\alpha$ in the simple roots, then $\alpha^\jepPubCoroot=\sum_{\beta\in\Pi}n_\beta(\alpha)\frac{\langle\beta,\alpha^\jepPubCoroot\rangle}{\langle\alpha,\beta^\jepPubCoroot\rangle}\beta^\jepPubCoroot$. If $\beta\in\Pi$ is a simple root, $\varpi_\beta$ denotes the fundamental weight corresponding to $\beta$. The set of simple roots $\Pi$ is a basis of $\mathfrak t^{\star}$, $\Pi^\jepPubCoroot=\{\beta^\jepPubCoroot\mid\beta\in\Pi\}$ is a basis of $\mathfrak t$ and $\{\varpi_\beta\mid\beta\in\Pi\}$ is by definition the basis of $\mathfrak t^{\star}$ dual to $\Pi^\jepPubCoroot$.

In this paper, we use the convention that if the root system $R$ of $\mathfrak g$ is simply laced (equivalently, of type A, D or E), then all the roots are long. Therefore short roots exist only if $R$ is not simply laced. There can be at most two root lengths in $R$. Moreover if there are two root lengths the ratio $\frac{\text{long root length}}{\text{short root length}}$ equals $\sqrt2$ because Hermitian symmetric spaces exist only when the type of $G$ is $A$, $B$, $C$, $D$, $E_6$ or $E_7$ (in particular, $G_2$ does not appear). We recall that for $\alpha,\beta\in R$, the ratio $\langle\alpha,\beta^\jepPubCoroot\rangle/\langle\beta,\alpha^\jepPubCoroot\rangle$ equals $(\text{root length of }\alpha)^2/(\text{root length of }\beta)^2$.

Linearly independent positive roots $\delta_1,\ldots,\delta_s$ are said to be \emph{strongly orthogonal} if for all $i\neq j$, $\delta_i\pm\delta_j$ is not a root. By \cite{jep93HC56}, or \cite[Ch.~VIII, \S7]{jep93Hel01}, we have
\begin{fact}
All maximal sets of noncompact strongly orthogonal roots have cardinality $r_M$.
\end{fact}
If $(\delta_1,\ldots,\delta_{r_M})$ is such a set then $\mathfrak g_{-\delta_1}\oplus\cdots\oplus\mathfrak g_{-\delta_{r_M}}$ is the tangent space to a maximal polydisc through the origin in $M$.

